\documentclass[fontsize=15pt]{jlreq}
%プリアンブル
\usepackage{amsmath}
\usepackage{mathtools}
\pagestyle{empty}
\usepackage[margin=10mm]{geometry}
\usepackage{tikz}
\usetikzlibrary{arrows.meta, calc}
\usepackage{here}

\newcommand{\m}[1]{\raisebox{0.1em}{\textcircled{\scriptsize #1}}}

\newcommand{\drawpyramid}[3]{%
  \begin{tikzpicture}[scale=0.7]
    \def\h{1.2} % 各段の高さ
    
    % 下段：植物（生産者）
    \draw[fill=gray!10, thick] (-#3/2, 0) rectangle (#3/2, \h);
    \node at (0, 0.5*\h) {\footnotesize 植物};
    
    % 中段：草食動物（一次消費者）
    \draw[fill=gray!25, thick] (-#2/2, \h) rectangle (#2/2, 2*\h);
    \node at (0, 1.5*\h) {\footnotesize 草食動物};
    
    % 上段：肉食動物（二次消費者）
    \draw[fill=gray!40, thick] (-#1/2, 2*\h) rectangle (#1/2, 3*\h);
    \node at (0, 2.5*\h) {\footnotesize 肉食動物};
  \end{tikzpicture}%
}



%本文
\begin{document}
\begin{center}
{\LARGE {振り返り}} \\
-生態- %単元ごとに変えるか消す
\end{center}
\vspace{1em}
\begin{flushright}
名前：\underline{\hspace{5cm}}
\end{flushright}

% 問題部分
\begin{enumerate}
  \item ある地域に生息，生育するすべての生物と，
  その地域の水や空気，土などの環境とをひとつのまとまりでとらえた
  ものを何というか．\\
  (\hspace{3cm})

  \vspace{2em}

  \item 食物連鎖とはなにか説明せよ．\\
  (\hspace{15cm})

  \vspace{2em}

  \item 一般的な生物の数量を多い順にかけ．(肉食動物，植物，草食動物から選ぶこと)\\
  (\hspace{4cm}$\to$\hspace{4cm}$\to$\hspace{4cm})

  \vspace{2em}

  \item 次の図は，様々な状態の
  生物ピラミッドを表している．今，
  図アのように何等かの原因によって植物が
  増えた状態となった．
  この後の変化として正しい順番になるように
  イ～エを並べ替えよ．

  \begin{figure}[H]
  \centering
  
  % --- 図1（ア）：植物だけが増加 ---1
  \begin{minipage}[t]{0.48\linewidth}
    \centering
    \drawpyramid{1.6}{3.0}{6.0}\\[0.5em]
    ア
  \end{minipage}
  \hfill
  %元 4
  \begin{minipage}[t]{0.48\linewidth}
    \centering
    \drawpyramid{1.6}{3.0}{4.4}\\[0.5em]
    イ
  \end{minipage}

  \vspace{2em} % 上下の間隔

  % --- 図3（ウ）：草食・肉食減少／植物回復 ---3
  \begin{minipage}[t]{0.48\linewidth}
    \centering
    \drawpyramid{1.2}{1.8}{4.2}\\[0.5em]
    ウ
  \end{minipage}
  \hfill
  %肉，草増 2
  \begin{minipage}[t]{0.48\linewidth}
    \centering
    \drawpyramid{2.0}{3.6}{4.0}\\[0.5em]
    エ
  \end{minipage}
\end{figure}
(ア$\to$\hspace{2cm}$\to$\hspace{2cm}$\to$\hspace{2cm})

\end{enumerate}

\end{document}